\section[Proposed Device Architecture]{Proposed Device Architecture and Design Rationale}
\label{sec:pd_design}

The preceding sections established the physical foundations
(\cref{sec:pd_fundamentals}), the performance metric language
(\cref{sec:pd_metrics}), and the architectural and material choices
(\cref{sec:pd_architectures,sec:pd_material}) that narrow the design space to
a vertical waveguide-integrated Ge-on-Si PIN photodetector at
\SI{1310}{\nano\metre}. This section specifies the device geometry and justifies
each design parameter in terms of the physics and metrics developed earlier.

\subsection[Vertical PIN Structure]{Vertical PIN Structure: Geometry, Layer Stack, and Field Distribution}
\label{subsec:pd_vertical_pin}

The photodiode is realised on a \SI{220}{\nano\metre} SOI platform. Each layer
parameter in \cref{tab:layer_stack} is the result of a specific physical argument
derived from the metric framework of \cref{sec:pd_metrics}.

\begin{table}[H]
\centering
\caption{Layer stack and primary material parameters of the Ge-on-Si n-i-p
photodiode simulation model. All doping concentrations are net values.}
\label{tab:layer_stack}
\footnotesize
\begin{tabular*}{\linewidth}{@{\extracolsep{\fill}}L{1.8cm}L{1.5cm}C{1.8cm}C{1.3cm}C{2.3cm}L{4.1cm}@{}}
\toprule
Layer & Material & Thickness & Type &
Concentration [\si{\per\centi\metre\cubed}] &
Primary function \\
\midrule
Cathode  & $n^{++}$-Ge & \SI{50}{\nano\metre}  & n & \SI{e19}{} &
    Ohmic contact; minimise FCA \\
Absorber & $i$-Ge      & \SI{350}{\nano\metre} & intrinsic & $\sim\SI{e13}{}$ &
    Optical absorption; depletion \\
Platform & $p^{++}$-Si & \SI{220}{\nano\metre} & p & \SI{e20}{} &
    Anode contact; current spreading \\
BOX      & \ce{SiO2}   & \SI{2}{\micro\metre}  & N/A & N/A &
    Optical isolation \\
Handle   & Si          & $\sim\SI{700}{\micro\metre}$ & N/A & N/A &
    Mechanical support \\
\bottomrule
\end{tabular*}
\end{table}

The intrinsic germanium thickness $W_i=\SI{350}{\nano\metre}$ is selected from
the optimisation of \cref{eq:wi_optimal}: it is within \SI{11}{\percent} of
$W_i^*\approx\SI{310}{\nano\metre}$ and delivers the conservative lower-bound
estimate $f_{\mathrm{tt}}\approx\SI{51}{\giga\hertz}$ from \cref{eq:bw_tt},
while the benchmark-calibrated two-carrier and peaking-enhanced response of the
same \SI{350}{\nano\metre} stack extends into the \SIrange{100}{120}{\giga\hertz}
class through \cref{eq:H_pd_pk}. The $n^{++}$-Ge cathode is limited to
\SI{50}{\nano\metre} to minimise free-carrier absorption (FCA) in the contact
region; this design choice simultaneously suppresses the diffusion tail in the
impulse response that would otherwise reduce the effective $f_{3\,\mathrm{dB}}$
below the prediction of \cref{eq:bw_tt}~\cite{Yan2024}.

\paragraph{Doping polarity and its impact on dark current.}
\label{subsec:pd_polarity}
The $p^{++}$-Si/$i$-Ge/$n^{++}$-Ge polarity (Si platform $p$-doped) places the
graded Ge/Si transition such that the valence-band discontinuity $\Delta E_v$
acts as an additive barrier to interface hole generation. In the alternative
$p$-on-$n$ configuration (Si platform $n$-doped), the same $\Delta E_v$ raises
$J_{\mathrm{surf}}$ above the value predicted by \cref{eq:jdark_surf} by modifying
the effective barrier for thermally generated holes. For a fixed reverse bias of
\SI{-1}{\volt}, the $p$-on-$n$ device exhibits $I_d$ a factor of approximately
two to three lower than the $n$-on-$p$ counterpart, translating through
\cref{eq:sens_dark} into a sensitivity improvement of approximately
\SI{1.5}{\decibel} in the dark-current-dominated regime~\cite{Alasio2024}.
The present device adopts the $p^{++}$-Si platform polarity precisely to
minimise the interface generation contribution while maintaining CMOS process
compatibility.

\begin{figure}[H]
    \centering
    \missingfigure[figwidth=0.86\linewidth]{
        Two-panel schematic of the selected Ge-on-Si PIN photodetector.
        Left panel: device cross-section used in the electrical model, showing
        the \SI{220}{\nano\metre} Si platform, the \SI{350}{\nano\metre}
        intrinsic Ge absorber, the doped contact regions, and the BOX.
        Right panel: top-view FDTD geometry showing the \SI{40}{\micro\metre}
        taper, the \SI{8}{\micro\metre}-long Ge absorber, the
        \SI{5}{\micro\metre} mesa width, the \SI{3}{\micro\metre} intrinsic
        region width, the U-shaped anode clearance, and the key stack
        specifications used in the optical and electrical models.
    }
    \caption{Selected device geometry for the Ge-on-Si PIN photodetector,
    showing the principal dimensions and model specifications used in the FDTD
    and CHARGE simulations.}
    \label{fig:design_dimensions}
\end{figure}

\subsection[Optical Coupling via Adiabatic Taper]{Optical Coupling via Adiabatic Taper}
\label{subsec:pd_taper}

Optical coupling from the input silicon waveguide is provided by a
\SI{40}{\micro\metre} adiabatic taper designed for near-unity coupling
efficiency. The taper adiabatically transforms the mode of the Si access
waveguide into the Ge-loaded mode over a length $L_t$ that is long enough to
suppress excitation of higher-order modes. By narrowing the Si waveguide tip to
approximately \SIrange{100}{200}{\nano\metre}, the guided mode is pushed
evanescently upward into the overlying Ge absorber. At the distal end of the
taper the mode is fully confined within the Ge, and absorption proceeds with
near-unity coupling efficiency. Evanescent coupling spreads the power transfer
over the full taper length, lowering the peak generation density for the same
absorbed power and preserving the geometric orthogonality of $L$ and
$W_i$~\cite{Wang2024,Yan2024}.

\subsection[U-Shaped Electrode Topology]{U-Shaped Electrode Topology: RC Reduction and Current Path Engineering}
\label{subsec:pd_electrode}

The U-shaped anode electrode is the primary design intervention for the RC
bandwidth in the vertical architecture. The anode electrode layout is the
dominant determinant of the RC bandwidth: selecting the wrong topology incurs
a penalty of approximately \SI{20}{\giga\hertz} in $f_{3\,\mathrm{dB}}$ relative
to the present state of the art.

\subsubsection{The Vertical n-i-p Equivalent Circuit}
\label{subsubsec:pd_eqckt}

The complete small-signal equivalent circuit of the waveguide-coupled vertical
n-i-p photodiode, including all parasitic elements, is shown in
\cref{fig:eq_circuit}. The photocurrent source $I_{\mathrm{ph}}=\mathcal{R}P_{\mathrm{opt}}$
drives the junction admittance, which consists of the junction capacitance $C_j$
in parallel with the shunt conductance $G_{\mathrm{sh}}=I_d/V_{\mathrm{bias}}$.
The junction is accessed through the series resistance network comprising the
cathode series resistance $R_{s,\mathrm{cat}}$ and the anode series resistance
$R_{s,\mathrm{an}}$ (path through the $p^{++}$-Si platform from the mesa base to
the metal anode pad). The bondpad-to-chip parasitic capacitance $C_{\mathrm{par}}$
shunts the junction at high frequencies, and the load termination
$R_L=\SI{50}{\ohm}$ converts the photocurrent into a voltage signal.

Lumping $R_s=R_{s,\mathrm{cat}}+R_{s,\mathrm{an}}$ and
$C_{\mathrm{tot}}=C_j+C_{\mathrm{par}}$, the voltage transfer function from
photocurrent source to load voltage is:
\begin{equation}
    H_{\mathrm{RC}}(f)
    = \frac{V_{\mathrm{out}}(f)}{I_{\mathrm{ph}}(f)}
    = \frac{R_L}{(1+j2\pi f R_L C_{\mathrm{tot}})(1+j2\pi f R_s C_{\mathrm{tot}})},
    \label{eq:hrc_full}
\end{equation}
which for $R_s\ll R_L$ reduces to the single-pole expression of \cref{eq:bw_rc}.

\begin{figure}[H]
\centering
\missingfigure[figwidth=0.84\linewidth]{
    Complete small-signal equivalent circuit of the vertical n-i-p detector,
    including the photocurrent source, junction capacitance, shunt
    conductance, parasitic capacitance, \SI{50}{\ohm} load, and the dominant
    anode resistance term whose value changes between the straight and
    U-shaped layouts.}
\caption{Complete small-signal equivalent circuit of the vertical n-i-p waveguide
photodiode including all parasitic elements from \cref{eq:hrc_full}. The dominant
bandwidth-limiting resistance is $R_{s,\mathrm{an}}$, the lateral conduction path
through the $p^{++}$-Si platform, which is reduced by \SI{36}{\percent} in the
U-shaped electrode topology relative to the parallel-electrode design.}
\label{fig:eq_circuit}
\end{figure}

\subsubsection{Parallel-Electrode versus U-shaped Anode: Quantitative Comparison}
\label{subsec:pd_electrode_compare}

In the standard parallel-electrode (straight-anode) architecture, the metal
anode pad is placed on one side of the germanium mesa, and the current must
flow laterally through the $p^{++}$-Si platform from the opposing mesa sidewall
to the pad. The mean lateral path length is approximately
$\ell_a^{\mathrm{par}}\approx W_m/2+d_{\mathrm{gap}}$, where $W_m$ is the mesa
width and $d_{\mathrm{gap}}$ is the clearance between the mesa sidewall and the
anode metal edge. For a \SI{5}{\micro\metre}-wide mesa and a \SI{0.5}{\micro\metre}
lithographic clearance:
\begin{equation}
    R_{s,\mathrm{an}}^{\mathrm{par}}
    \approx \frac{\rho_{\mathrm{Si}}}{t_{\mathrm{Si}}}
    \cdot\frac{\ell_a^{\mathrm{par}}}{L}
    = \frac{\rho_{\mathrm{Si}}}{t_{\mathrm{Si}}}
    \cdot\frac{W_m/2+d_{\mathrm{gap}}}{L}.
    \label{eq:rs_parallel}
\end{equation}
In the U-shaped anode architecture, the metal wraps around three sides of the
mesa (the two long sidewalls and one short end), so that current from any point
on the mesa base must travel at most $W_m/2$ laterally before reaching a
metal contact. The effective mean path length is reduced to~\cite{Shi2024}:
\begin{equation}
    \ell_a^{\mathrm{U}} \approx \frac{W_m}{4} + d_{\mathrm{gap}},
    \label{eq:la_u}
\end{equation}
giving a series resistance:
\begin{equation}
    R_{s,\mathrm{an}}^{\mathrm{U}}
    \approx \frac{\rho_{\mathrm{Si}}}{t_{\mathrm{Si}}}
    \cdot\frac{W_m/4+d_{\mathrm{gap}}}{L}.
    \label{eq:rs_u}
\end{equation}
The reduction ratio is:
\begin{equation}
    \frac{R_{s,\mathrm{an}}^{\mathrm{U}}}{R_{s,\mathrm{an}}^{\mathrm{par}}}
    = \frac{W_m/4+d_{\mathrm{gap}}}{W_m/2+d_{\mathrm{gap}}}
    \approx 0.60 \quad\text{for } W_m=\SI{5}{\micro\metre},\;d_{\mathrm{gap}}=\SI{0.5}{\micro\metre},
    \label{eq:rs_ratio}
\end{equation}
corresponding to a \SI{40}{\percent} resistance reduction, consistent with the
\SI{36}{\percent} reported experimentally by Shi~\textit{et al.}~\cite{Shi2024}.
Substituting into \cref{eq:bw_rc}:
\begin{equation}
    \frac{f_{\mathrm{RC}}^{\mathrm{U}}}{f_{\mathrm{RC}}^{\mathrm{par}}}
    = \frac{R_{\mathrm{tot}}^{\mathrm{par}}}{R_{\mathrm{tot}}^{\mathrm{U}}}
    = \frac{R_L + R_{s,\mathrm{an}}^{\mathrm{par}}}{R_L + R_{s,\mathrm{an}}^{\mathrm{U}}}.
    \label{eq:frc_ratio}
\end{equation}
For $R_L=\SI{50}{\ohm}$, $R_{s,\mathrm{an}}^{\mathrm{par}}=\SI{8}{\ohm}$, and
$R_{s,\mathrm{an}}^{\mathrm{U}}=\SI{5}{\ohm}$, this gives
$f_{\mathrm{RC}}^{\mathrm{U}}/f_{\mathrm{RC}}^{\mathrm{par}}=58/55\approx1.055$,
a modest \SI{5.5}{\percent} increase in $f_{\mathrm{RC}}$ alone. The larger
bandwidth gain emerges only when the electrode is treated as part of the full
electro-optical co-design problem. In the benchmark-calibrated peaked circuit of
\cref{eq:H_pd_pk}, the reduced $R_s$ also enables a shallower and more useful
inductive peaking condition, and the experimentally observed bandwidth increases
from \SI{83}{\giga\hertz} (parallel electrode) to
\SI{103}{\giga\hertz} (single-U)~\cite{Shi2024}.

\begin{figure}[H]
\centering
\missingfigure[figwidth=0.86\linewidth]{
    Four-panel topology figure. Straight and U-shaped electrode layouts are
    compared in plan view, their current paths are mapped to the dominant
    anode resistance, and the resulting bandwidth shift is shown relative to
    the lower-bound transport ceiling and the benchmark-calibrated peaked
    response.}
\caption{Parallel-electrode versus U-shaped anode comparison. The current path
in the $p^{++}$-Si platform is halved by the U-shape (left to centre panels),
reducing $R_{s,\mathrm{an}}$ by \SI{36}{\percent} from \cref{eq:rs_ratio}.}
\label{fig:electrode_compare}
\end{figure}

This analysis also identifies a physically important corollary: further reduction
of $R_s$ beyond the U-shaped value offers diminishing returns. In the
simplified screening model of \cref{eq:bw_total}, letting $R_s\to0$ drives the
response toward the lower-bound transport limit of \cref{eq:bw_tt}; in the
benchmark-calibrated full-response picture of \cref{eq:H_pd_pk}, the practical
ceiling of the \SI{350}{\nano\metre} vertical stack lies in the
\SIrange{110}{120}{\giga\hertz} class rather than increasing indefinitely with
electrode optimisation.

\paragraph{Mask-level representation in KLayout.}
The simulated device is tied to an explicit lithographic layout through the
KLayout script \texttt{system-level/ge\_pd\_layout.py}, which writes the cell
\texttt{GE\_PD\_USHAPED\_OBAND} to \texttt{ge\_pd\_layout.gds}. The layout uses
a database unit of \SI{1}{\nano\metre} and maps the principal process features
onto ten logical layers corresponding to the Si slab, Si ridge, Si taper, Ge,
$p^{++}$ region, $n^{++}$ region, via level, M1, M2, and BOX.

\begin{figure}[H]
    \centering
    \missingfigure[figwidth=0.86\linewidth]{
        Top-view KLayout rendering of the mask-level photodetector layout,
        showing the Si taper, Ge mesa, U-shaped M1 anode, via array, and the
        M2 runner derived from the scripted geometry in
        \texttt{system-level/ge\_pd\_layout.py}.
    }
    \caption{Mask-level layout of the simulated O-band photodetector derived
    from the KLayout script. The layout preserves the optical path, contact
    topology, and routing footprint used to define the device implementation.}
    \label{fig:klayout_layout}
\end{figure}

The geometry parameters defined above, namely $W_i$, $L$, $W_m$, the electrode
topology, and the taper length, constitute the degrees of freedom of the design
space. The quantitative mapping from these parameters to the performance metrics
of \cref{sec:pd_metrics}, including the absorber length sweep, the
intrinsic-thickness trade-off, and the electric-field screening threshold, is
presented with simulation data in \cref{sec:pd_results}.
